\chapter{Introducción}
\label{ch:introduccion}

\section{Propósito y alcance}

El presente documento corresponde a la \textbf{Revisión Crítica de Diseño (CDR)}
del \textit{CMOS Microscopy Payload} (CMOSM) del CubeSat INTISAT, desarrollado en
el Instituto de Radioastronomía (INRAS) de la Pontificia Universidad Católica del
Perú (PUCP).

El objetivo principal del payload CMOSM es demostrar un sistema de
\textbf{microscopía sin lentes (lensless) en configuración de contacto}, basado
en un sensor CMOS comercial, capaz de capturar imágenes de muestras biológicas
planas con suficiente contraste y resolución para la identificación de
estructuras celulares, tales como paredes celulares en tejidos vegetales.

El sistema prioriza la obtención de imágenes reproducibles mediante un control
adecuado de la iluminación, el registro de parámetros de captura y el correcto
\textbf{acoplamiento óptico entre la muestra y el sensor}. Este último se
identifica como un factor crítico para maximizar el contraste y la definición
de bordes en microscopía de contacto, y es validado experimentalmente en la
fase actual del proyecto.
\section{Alcance técnico del payload}

El diseño del CMOSM se fundamenta en un enfoque de microscopía de contacto, en el
cual la muestra se coloca a una distancia del orden de micrómetros del plano del
sensor CMOS. En este régimen, la formación de imagen está dominada por fenómenos
de absorción, dispersión cercana y sombreado, siendo altamente sensible a la
interfaz óptica muestra--sensor.

Durante la fase actual de desarrollo, la interfaz de muestra se implementa
mediante soluciones experimentales de laboratorio, empleadas con el objetivo de
validar el principio físico de acoplamiento óptico y su impacto en la calidad de
imagen. La definición de un material definitivo para dicha interfaz se
identifica como trabajo futuro y queda fuera del alcance del presente CDR.

Sobre esta base física, el payload incorpora capacidades de procesamiento
computacional orientadas a la mejora de la legibilidad visual de las imágenes
adquiridas. Técnicas avanzadas como la Fourier Ptychographic Microscopy (FPM) y
la super--resolución basada en aprendizaje profundo se consideran
\textbf{modos complementarios y experimentales}, evaluados en función de su
aporte incremental sobre la imagen base capturada directamente por el sistema.
Las capacidades principales del sistema incluyen:

\begin{itemize}
    \item Captura directa de imágenes mediante un sensor CMOS comercial
    operando en configuración sin lentes.
    \item Control de iluminación estructurada y de intensidad para optimizar el
    contraste de la muestra.
    \item Registro completo de parámetros de captura (exposición, ganancia,
    configuración de iluminación y métricas de intensidad) como metadatos
    asociados a cada imagen.
    \item Procesamiento computacional posterior para realce visual, incluyendo
    técnicas de super--resolución.
    \item Evaluación experimental de técnicas de reconstrucción avanzadas,
    como Fourier Ptychographic Microscopy (FPM).
\end{itemize}
\section{Estructura del documento}

El documento se organiza siguiendo una estructura inspirada en los
\emph{Design Definition \& Justification Files} (DDJF) utilizados en misiones
CubeSat de referencia y alineada con lineamientos ECSS/NASA, priorizando la
descripción del sistema físico de adquisición, su validación experimental y la
justificación de las decisiones de diseño adoptadas.

\begin{itemize}
    \item \textbf{Capítulo~\ref{ch:introduccion}}: contexto general, propósito,
    documentos aplicables y definiciones.
    \item \textbf{Capítulo~\ref{ch:requerimientos}}: definición de
    requerimientos de misión y de payload.
    \item \textbf{Capítulo~\ref{ch:modos-operacion}}: modos de operación,
    secuencias de captura y procedimientos nominales y anómalos.
    \item \textbf{Capítulo~\ref{ch:funcional}}: descripción funcional de la
    arquitectura del payload y descomposición en subsistemas.
    \item \textbf{Capítulo~\ref{ch:diseno-justificacion}}: diseño detallado y
    justificación de las decisiones ópticas, mecánicas, electrónicas y de
    software.
    \item \textbf{Capítulo~\ref{ch:verificacion}}: estrategia de verificación,
    matriz de cumplimiento y resumen de pruebas.
    \item \textbf{Capítulo~\ref{ch:riesgos}}: identificación de riesgos
    técnicos y operacionales, con sus mitigaciones.
    \item \textbf{Capítulo~\ref{ch:conclusiones}}: conclusiones de la CDR y
    recomendaciones para la siguiente fase.
    \item \textbf{Apéndices}: información de soporte y resultados detallados.
\end{itemize}
\section{Contexto del proyecto INTISAT}

INTISAT es un proyecto CubeSat académico del INRAS orientado a la validación de
tecnologías de observación y experimentación científica en órbita baja (LEO).
Dentro de su arquitectura de cargas útiles, el payload CMOSM se concibe como un
\textbf{demostrador tecnológico} de microscopía sin lentes de contacto y
procesamiento computacional asociado, reutilizable en futuras misiones
educativas o científicas.

El diseño del CMOSM prioriza:
\begin{itemize}
    \item Uso de componentes \emph{Commercial Off-The-Shelf} (COTS) como la
    cámara CMOS OV5647.
    \item Arquitectura mecánica simple, ligera y replicable en laboratorio.
    \item Alto nivel de telemetría y trazabilidad de parámetros de captura para
    análisis en tierra.
    \item Modularidad del software para permitir iteración rápida sobre
    algoritmos de procesamiento y reconstrucción experimental.
\end{itemize}

\begin{figure}[H]
    \centering
    \includegraphics[width=0.6\textwidth]{figures/IntiSat_Renderization_CMOS.png}
    \caption{Renderizado del CubeSat INTISAT mostrando la configuración del payload CMOSM.}
    \label{fig:intisat_render}
\end{figure}
